% ============================================
% РАЗДЕЛ 1. АНАЛИЗ ПРЕДМЕТНОЙ ОБЛАСТИ
% ============================================

\section{Анализ предметной области}

\subsection{Задача распознавания ключевых слов для устройств интернета вещей}

Голосовое управление устройствами интернета вещей (Internet of Things, IoT) предполагает наличие механизма обнаружения заранее определённых слов-команд в непрерывном аудиопотоке. В литературе данная задача обозначается термином Keyword Spotting (KWS) и представляет собой частный случай распознавания речи с ограниченным словарём~--- от единиц до нескольких десятков команд~\cite{src:1}. В отличие от полноценных систем автоматического распознавания речи, обрабатывающих произвольные фразы, системы KWS ориентированы на минимизацию вычислительных ресурсов и энергопотребления за счёт упрощения модели, что делает их пригодными для развёртывания непосредственно на ресурсоограниченных встраиваемых платформах~\cite{src:1, src:8}.

Систематизация существующих алгоритмических решений для развёртывания нейронных сетей на устройствах с ограниченными ресурсами проведена в обзорной работе Khan и др.~\cite{src:27}. Авторы выделяют три взаимодополняющих направления адаптации моделей: оптимизация архитектуры (использование компактных топологий с уменьшенным числом параметров), посттренировочные преобразования (квантизация, обрезка соединений, дистилляция знаний) и аппаратно-программный кодизайн (специализированные ядра, эксплуатация SIMD-расширений). Указанная классификация подтверждает методологическую состоятельность подхода, принятого в настоящей работе: компактная архитектура DS-CNN, квантизация INT8 и развёртывание через библиотеку аппаратно-ускоренных ядер ESP-NN относятся к перечисленным направлениям и применяются совместно.

Практические сценарии применения KWS в IoT охватывают широкий класс задач: голосовая активация промышленных носимых устройств в условиях занятых рук оператора, речевое управление автономными роботизированными системами, активация охранных узлов по голосовой команде и управление бытовыми устройствами умного дома. Во всех перечисленных сценариях система должна находиться в состоянии постоянной готовности к приёму команды, что предполагает непрерывную работу аудиопайплайна: захват аудиосигнала, извлечение признаков и классификацию. При батарейном питании устройства данное требование вступает в прямое противоречие с ограниченным энергетическим бюджетом. Указанное противоречие между необходимостью непрерывного мониторинга и минимизацией среднего энергопотребления определяет актуальность разработки специализированных энергоэффективных архитектур KWS-систем.

В качестве стандартного бенчмарка для оценки и сравнения моделей KWS в научном сообществе используется открытый набор данных Google Speech Commands~\cite{src:2}. Набор версии~2 содержит около 105~000 однословных аудиозаписей длительностью в одну секунду, охватывающих 35 англоязычных команд, записанных более чем 2~600 дикторами в различных акустических условиях. В стандартной конфигурации из 35 команд выделяются 10 целевых слов (\textit{yes, no, up, down, left, right, on, off, stop, go}), отдельный класс тишины (silence) и класс неизвестных слов (unknown), формируя задачу классификации на 12 классов~\cite{src:2}. Именно данная конфигурация принята в бенчмарке MLPerf Tiny~\cite{src:5} в качестве эталонной для задачи KWS на микроконтроллерах.

\subsection{Предварительная обработка аудиосигнала: извлечение признаков}

Нейросетевые модели KWS, как правило, не обрабатывают сырой аудиосигнал непосредственно. Входным представлением служат спектрально-временные признаки, извлечённые из аудиопотока на этапе предобработки. Наиболее распространённым представлением являются мел-частотные кепстральные коэффициенты (Mel-Frequency Cepstral Coefficients, MFCC), предложенные Davis и Mermelstein в 1980 году~\cite{src:10}.

Процедура извлечения MFCC включает следующие этапы: сегментацию аудиосигнала на перекрывающиеся фреймы (типичная длина окна 25--40~мс, шаг 10--20~мс); применение оконной функции (Хэмминга или Ханна) для подавления спектральной утечки; вычисление дискретного преобразования Фурье (ДПФ) для получения амплитудного спектра каждого фрейма; применение набора треугольных мел-фильтров, равномерно распределённых по психоакустической мел-шкале; логарифмирование энергии в каждой мел-полосе; и, наконец, дискретное косинусное преобразование (ДКП), результатом которого являются декоррелированные кепстральные коэффициенты. Нижние коэффициенты описывают грубую спектральную огибающую звука~--- фактически формантную структуру речевого тракта,~--- а верхние фиксируют тонкие спектральные детали. Для задач KWS обычно сохраняются первые 10--13 коэффициентов из каждого фрейма.

В результате одна секунда аудиосигнала с частотой дискретизации 16~кГц преобразуется в матрицу фиксированного размера~--- например, $49 \times 10$ при параметрах, использованных в работе Zhang и др.~\cite{src:1}: длина окна 40~мс, шаг 20~мс, 10 MFCC-коэффициентов. Полученная матрица может быть интерпретирована как одноканальное <<изображение>>, где строки соответствуют временны\'{м} фреймам, а столбцы~--- спектральным характеристикам. Именно данная интерпретация обосновывает применение свёрточных нейронных сетей для классификации.

Выбор количества сохраняемых кепстральных коэффициентов представляет самостоятельный аспект проектирования. Полный набор из $M$ коэффициентов (по числу мел-фильтров) избыточен: верхние коэффициенты описывают высокочастотные детали спектра, слабо коррелирующие с фонемным содержанием. Сокращение размерности признакового пространства не только уменьшает входной тензор, но и снижает риск переобучения. В работе Исхаковой и др.~\cite{src:26} применительно к задаче распознавания речевых эмоций экспериментально показано, что снижение размерности признакового пространства с использованием свёрточной нейронной сети сохраняет основные классификационные свойства при существенном сокращении вычислительных затрат. Аналогичный принцип используется и в настоящей работе: из 40 энергий мел-фильтров после ДКП сохраняются только первые 10 коэффициентов, что соответствует общепринятой практике в задачах KWS.

Целесообразность совместного использования MFCC и свёрточных нейронных сетей подтверждается смежными задачами обработки речевого сигнала. Так, в работе Муртазина и др.~\cite{src:24} представлен алгоритм выявления синтезированного голоса, основанный на сочетании кепстральных коэффициентов и свёрточной нейронной сети; авторы продемонстрировали, что подобная двухступенчатая комбинация (психоакустически мотивированное преобразование признаков и последующая нейросетевая классификация) обеспечивает высокую точность при ограниченной вычислительной сложности модели. Применительно к KWS подобный подход также является доминирующим, что дополнительно обосновывает выбор пары MFCC+DS-CNN в настоящей работе. Отметим, что устойчивость к атакам синтезированной речью является самостоятельной проблемой безопасности голосовых интерфейсов IoT, выходящей за рамки настоящего исследования.

Значимость этапа предобработки с точки зрения энергопотребления подчёркивается в работе Cerutti и др.~\cite{src:7}: экспериментально установлено, что в стандартном KWS-пайплайне на микроконтроллере до 85\% совокупной энергии расходуется на захват аудиосигнала и извлечение MFCC-признаков, тогда как на инференс нейронной сети приходятся оставшиеся 15\%. Данный результат указывает на то, что оптимизация предобработки может оказать большее влияние на энергоэффективность, чем оптимизация самой модели классификатора.

\subsection{Классификация подходов к реализации KWS}

Существующие подходы к построению систем KWS для IoT-устройств могут быть классифицированы по месту выполнения инференса и типу вычислительной платформы.

\subsubsection{Облачные решения}

При облачном подходе устройство осуществляет захват аудиоданных и передачу их на удалённый сервер для инференса полноразмерной модели (Conformer, Whisper и др.). К достоинствам относятся неограниченная сложность моделей, высокая точность и поддержка произвольного словаря. Вместе с тем для автономных IoT-устройств облачный подход обладает существенными ограничениями. Передача аудиоданных по Wi-Fi требует поддержания активного радиоканала: согласно технической документации ESP32-S3~\cite{src:9}, потребление в режиме передачи Wi-Fi составляет 180--240~мА при 3,3~В, что соответствует мощности 600--800~мВт. Помимо энергопотребления радиомодуля, облачный подход привносит сетевую латентность порядка 50--500~мс, зависит от наличия и качества сетевого подключения и ставит под вопрос приватность голосовых данных пользователя. По совокупности факторов облачный подход наименее пригоден для батарейных автономных устройств.

\subsubsection{Специализированные аппаратные платформы}

\textit{ASIC (Application-Specific Integrated Circuit).} Процессоры, разработанные исключительно для нейросетевого инференса, обеспечивают наивысшую энергоэффективность среди рассмотренных решений. Линейка Syntiant NDP100/NDP101/NDP120 реализует архитектуру всегда активного прослушивания на базе специализированного вычислительного ядра Syntiant Core~2. По данным бенчмарка MLPerf Tiny~\cite{src:5}, данные процессоры демонстрируют потребление порядка единиц микроджоулей на один акт инференса~--- на два-три порядка ниже, чем MCU общего назначения. Однако закрытая архитектура, невозможность модификации модели или аппаратной логики, высокая стоимость единичных образцов и ограниченная доступность на российском рынке делают ASIC неприменимыми в условиях данного исследования.

\textit{FPGA (Field-Programmable Gate Array).} Программируемые логические интегральные схемы позволяют реализовать аппаратные ускорители нейросетевых вычислений с полным контролем над микроархитектурой. В рамках бенчмарка MLPerf Tiny Duarte и др.~\cite{src:21} продемонстрировали развёртывание KWS-модели на FPGA-платформах Xilinx Pynq-Z2 и Arty~A7-100T с использованием открытых фреймворков hls4ml и FINN, достигнув латентности менее 20~мкс и энергопотребления от 30~мкДж на один инференс. Данные результаты свидетельствуют о потенциале FPGA для задач с ультранизкой латентностью. Вместе с тем разработка на FPGA предполагает проектирование цифровых схем на языках описания аппаратуры (VHDL, Verilog), что требует специализированных компетенций, не входящих в стандартный стек разработки IoT-устройств. Сложность данной задачи иллюстрируется работой Ушениной и Данилова~\cite{src:25}: проектирование на ПЛИС даже отдельного модуля искусственного нейрона для последовательно-параллельных архитектур нейронных сетей с прямой связью представляет самостоятельную инженерную задачу, требующую тщательной оптимизации последовательности операций умножения с накоплением и распределения ресурсов между арифметическими блоками и блочной памятью. При переходе от единичного нейрона к полной свёрточной сети уровня DS-CNN сложность проектирования возрастает на порядки. Кроме того, FPGA не содержат встроенной периферии для IoT-приложений (контроллеры Wi-Fi, BLE, I2S), что вынуждает использовать внешние компоненты и увеличивает сложность конечной системы. Стоимость отладочных плат FPGA с достаточным объёмом логических ресурсов (от десятков долларов для Lattice iCE40 UltraPlus до сотен для Xilinx Zynq) также существенно превышает стоимость микроконтроллерных платформ.

\textit{DSP (Digital Signal Processor).} Цифровые сигнальные процессоры оптимизированы для потоковых вычислений: свёрток, быстрого преобразования Фурье и матричных операций. Архитектуры семейства Texas Instruments TMS320C5x обеспечивают потребление порядка единиц милливатт при выполнении аудиообработки. Однако DSP ориентированы на классические алгоритмы обработки сигналов, а не на произвольные нейросетевые архитектуры; не обладают периферией, характерной для IoT-платформ; и, как правило, требуют дополнительных MCU для реализации управляющей логики. В контексте KWS использование выделенного DSP целесообразно главным образом в высокопроизводительных системах с непрерывной обработкой нескольких аудиоканалов, что выходит за рамки типичных IoT-сценариев.

\subsubsection{Микроконтроллеры общего назначения}

Микроконтроллеры (MCU) с поддержкой целочисленных SIMD-расширений и аппаратных режимов энергосбережения представляют третий подход к реализации KWS на периферии сети. Инференс квантизированной нейросетевой модели выполняется непосредственно на ЦПУ MCU, предобработка~--- там же.

Среди платформ, применяемых для KWS, выделяются:

\begin{itemize}
	\item семейство ARM Cortex-M (STM32L4, STM32U3, STM32U5, STM32H7, nRF52840) с набором SIMD-инструкций DSP Extension и библиотекой CMSIS-NN для аппаратного ускорения целочисленных нейросетевых операций; в рамках бенчмарка MLPerf Tiny v1.3 плата NUCLEO-U385RG-Q на базе STM32U3 продемонстрировала 48 инференсов в секунду при потреблении 245~мВт для задачи KWS;
	\item платформы Xtensa (Espressif ESP32-S3) с двухъядерным процессором LX7, 128-битными SIMD-регистрами, встроенным контроллером I2S для цифровых MEMS-микрофонов, модулями Wi-Fi и BLE, а также аппаратным режимом глубокого сна с пробуждением по RTC GPIO~\cite{src:9};
	\item платформы RISC-V (Espressif ESP32-C3, GreenWaves GAP8/GAP9) с открытым набором инструкций; GAP8 интегрирует кластер из 8+1 ядер RISC-V с аппаратным конвейером для параллельных MAC-операций, занимая промежуточное положение между MCU и ASIC по энергоэффективности.
\end{itemize}

Типичное потребление MCU при активном инференсе KWS составляет 70--150~мВт~\cite{src:8}, что на порядки выше ASIC. Ключевым методом компенсации данного разрыва является каскадная архитектура: MCU находится в глубоком сне (единицы мкА), просыпается по внешнему событию и возвращается в сон после обработки. Высокое мгновенное потребление компенсируется тем, что активный режим занимает малую долю общего времени.

\subsection{Каскадные подходы к снижению энергопотребления KWS}

Каскадная детекция разделяет аудиопайплайн на несколько последовательных этапов с возрастающей вычислительной сложностью: каждый этап выступает фильтром-предусловием для активации следующего. В периоды тишины работает только первый, наиболее дешёвый этап; нейросетевой инференс запускается лишь при обнаружении потенциального речевого события.

Ключевой компонент каскадной архитектуры~--- детектор голосовой (или речевой) активности (Voice Activity Detector, VAD)~--- является самостоятельным направлением исследований в области цифровой обработки сигналов. В отечественной литературе вопросы построения детекторов VAD рассмотрены, в частности, в работе Стефаниди и др.~\cite{src:23}, предложивших комбинированный детектор, объединяющий несколько признаков речевой активности: энергию сигнала, частоту пересечения нуля и спектральные характеристики. Такая комбинация обеспечивает устойчивость к разнотипным акустическим помехам по сравнению с одноэлементными пороговыми решениями. Предельным упрощением VAD-схемы является аналоговый пороговый компаратор, фиксирующий лишь превышение огибающей сигнала заданного уровня. Подобный детектор не отделяет речь от прочих громких звуков, но при потреблении в сотни микроватт служит эффективным <<нулевым>> фильтром каскада, отсекающим длительные периоды абсолютной тишины. Именно такая упрощённая аналоговая реализация применяется в настоящей работе как первый этап каскада; перенос алгоритмически сложных комбинированных VAD на маломощный сопроцессор ULP составляет самостоятельное направление развития предложенной системы.

Hardy и Badets~\cite{src:6} предложили двухуровневый каскад: аналоговый детектор голосовой активности с потреблением порядка 20~мкВт управляет пробуждением основного MCU для запуска RNN-классификатора. Данный подход продемонстрировал устойчивость к реальным акустическим условиям при потреблении, сопоставимом со специализированными ASIC.

Cerutti и др.~\cite{src:7} реализовали систему с аналоговым извлечением бинарных признаков и бинарной нейронной сетью на MCU, достигнув общего потребления менее 1~мВт в режиме непрерывного мониторинга. Замена цифровой предобработки аналоговым фронтендом позволила сократить энергозатраты на этап захвата и извлечения признаков в 29 раз, снизив их долю с доминирующих 85\% до 16\% общего энергобюджета. Данный результат экспериментально подтверждает, что основным источником энергопотребления в KWS-пайплайне на MCU является не инференс нейронной сети, а предшествующая ему обработка аудиосигнала.

Bushur и Chen~\cite{src:8} провели сравнительное исследование нейросетевых архитектур для KWS на различных периферийных платформах, оценивая компромисс между точностью, латентностью и размером модели. Их работа акцентирует внимание на необходимости учёта аппаратных особенностей целевой платформы при выборе и оптимизации модели.

Анализ перечисленных работ выявляет характерную особенность: каждая из них рассматривает лишь отдельный аспект проблемы энергоэффективности KWS. Hardy и Badets исследуют каскадную архитектуру, но не применяют квантизацию модели и не проводят декомпозицию энергобюджета по компонентам. Cerutti и др. оптимизируют предобработку, но не используют каскадное пробуждение из режима глубокого сна с многоступенчатой фильтрацией событий. Bushur и Chen сравнивают платформы и архитектуры, но не рассматривают каскадный подход. Комплексное решение, объединяющее каскадную архитектуру с квантизацией модели и количественной декомпозицией энергобюджета по этапам обработки, в рассмотренной литературе представлено недостаточно. Данное наблюдение определяет позиционирование настоящей работы.

\subsection{Архитектуры нейронных сетей для распознавания ключевых слов}

Развитие нейросетевых подходов к задаче KWS прошло путь от полносвязных сетей (DNN) через свёрточные (CNN) к специализированным компактным архитектурам, оптимизированным для встраиваемых платформ. Ниже рассматриваются архитектуры, наиболее значимые для задачи KWS на микроконтроллерах.

\textit{DS-CNN (Depthwise Separable CNN).} Архитектура, предложенная Zhang и др.~\cite{src:1} в работе <<Hello Edge>>, основана на замене стандартных свёрточных слоёв двухступенчатой декомпозицией: пространственная свёртка по каждому каналу независимо (depthwise convolution) и линейная комбинация каналов с ядром $1 \times 1$ (pointwise convolution). Данный приём, заимствованный из архитектуры MobileNet~\cite{src:11}, сокращает количество параметров и MAC-операций (умножений с накоплением) в 8--10 раз при сопоставимой точности. На наборе Google Speech Commands архитектура DS-CNN достигает точности 95,4\%, что на 10 процентных пунктов превышает показатели полносвязной сети DNN с аналогичным количеством параметров~\cite{src:1}.

В работе~\cite{src:1} предложены три варианта масштабирования, привязанные к ресурсным ограничениям целевых платформ:

\begin{itemize}
	\item DS-CNN-S (Small): ограничение 80~КБ памяти, 6 миллионов MAC-операций;
	\item DS-CNN-M (Medium): ограничение 200~КБ памяти, 20 миллионов MAC-операций;
	\item DS-CNN-L (Large): ограничение 500~КБ памяти, 80 миллионов MAC-операций.
\end{itemize}

В бенчмарке MLPerf Tiny~\cite{src:5} архитектура DS-CNN принята как эталонная для задачи KWS~--- все участники бенчмарка обязаны обеспечить точность не ниже 90\% на наборе Google Speech Commands.

\textit{TC-ResNet (Temporal Convolution ResNet).} Архитектура, предложенная Choi и др.~\cite{src:12}, применяет одномерные свёртки вдоль временной оси с остаточными соединениями (residual connections). В отличие от DS-CNN, использующей двумерные свёртки по матрице MFCC-признаков, TC-ResNet обрабатывает MFCC-фреймы как одномерные временные ряды, что позволяет существенно сократить количество вычислений. Модель TC-ResNet8 продемонстрировала ускорение в 1,5--15 раз по сравнению с вариантами DS-CNN-S/M/L при повышении точности на 0,7--1,7 процентного пункта, а также ускорение в 385 раз по сравнению с глубокой моделью Res15 при сопоставимой точности (96,6\%)~\cite{src:12}. Вместе с тем архитектура менее исследована в контексте квантизации и развёртывания через TFLite Micro на платформах Xtensa.

\textit{CRNN (Convolutional Recurrent Neural Network).} Гибридная архитектура, комбинирующая свёрточные слои для извлечения локальных спектрально-временных признаков с рекуррентными слоями (GRU или LSTM) для моделирования долгосрочных временных зависимостей. Обеспечивает высокую точность, однако рекуррентные слои существенно увеличивают требования к оперативной памяти из-за необходимости хранения скрытого состояния между фреймами, что ограничивает применимость CRNN на MCU с типичным объёмом SRAM 256--520~КБ.

Для настоящей работы выбрана архитектура DS-CNN, что обосновано совокупностью факторов: наличие эталонной реализации в бенчмарке MLPerf Tiny~\cite{src:5}; доказанная эффективность на микроконтроллерных платформах; масштабируемость (варианты S/M/L позволяют исследовать компромисс между размером модели и точностью); эффективная утилизация SIMD-инструкций архитектуры Xtensa LX7~--- pointwise-свёртки (ядро $1 \times 1$) сводятся к матричным умножениям INT8-элементов, при этом в 128-битный SIMD-регистр помещаются 16 весов формата INT8, что обеспечивает 16 MAC-операций за одну инструкцию~\cite{src:9}.

\subsection{Методы квантизации нейронных сетей}

Квантизация нейронных сетей~--- процедура снижения разрядности представления весов и активаций модели~--- является ключевым методом адаптации нейросетевых моделей к ресурсным ограничениям микроконтроллеров. Типичным является переход от 32-битных чисел с плавающей точкой (FP32) к 8-битным целым числам (INT8)~\cite{src:3}. Применительно к задаче KWS квантизация обеспечивает несколько взаимосвязанных преимуществ.

Во-первых, сокращение размера модели: при переходе от FP32 к INT8 каждый вес занимает 1 байт вместо 4, что обеспечивает теоретический коэффициент сжатия до $4\times$. Для модели DS-CNN-M, занимающей около 956~КБ в формате FP32, размер после квантизации составляет порядка 260~КБ~--- существенно ниже объёма флеш-памяти типичных MCU (4--8~МБ).

Во-вторых, ускорение инференса: целочисленная арифметика INT8 нативно поддерживается SIMD-расширениями современных MCU. На платформе ESP32-S3 библиотека ESP-NN реализует оптимизированные ядра свёрточных операций с использованием 128-битных векторных регистров архитектуры Xtensa LX7~\cite{src:9}, обеспечивая кратное ускорение по сравнению с операциями с плавающей точкой.

В-третьих, снижение энергопотребления~--- как прямое (целочисленные операции потребляют меньше энергии), так и косвенное: сокращённое время инференса позволяет раньше вернуть MCU в режим глубокого сна, уменьшая среднее потребление в каскадной архитектуре.

В литературе выделяют два основных подхода к квантизации~\cite{src:3}.

\textit{Посттренировочная квантизация (Post-Training Quantization, PTQ)} выполняется после завершения обучения модели. Предобученная модель FP32 преобразуется к INT8-представлению на основании статистик распределения весов и активаций, собранных на репрезентативном калибровочном наборе данных. Для каждого тензора вычисляются два параметра~--- масштаб (scale) и нулевая точка (zero point),~--- определяющие линейное отображение из диапазона вещественных чисел в диапазон целых. Достоинством PTQ является простота: процедура не требует повторного обучения, выполняется за минуты и применима к произвольной предобученной архитектуре. Недостатком~--- потенциально большая потеря точности, поскольку модель не компенсирует ошибки квантования.

\textit{Квантизация с учётом при обучении (Quantization-Aware Training, QAT)} встраивает процедуру квантизации в цикл обучения. На этапе прямого прохода в граф вычислений вставляются узлы симуляции квантования (fake quantization nodes), моделирующие эффект округления активаций до INT8. Градиенты рассчитываются по приближению Straight-Through Estimator, что позволяет оптимизатору компенсировать ошибки округления. QAT обеспечивает меньшую деградацию точности по сравнению с PTQ, что особенно заметно при квантизации ниже 8 бит~\cite{src:3}. Недостатком является необходимость повторного обучения и более сложная инфраструктура тренировочного пайплайна.

Для задачи KWS с квантизацией до INT8 оба подхода обеспечивают потерю точности менее 1 процентного пункта~\cite{src:1, src:3}. Тем не менее их сравнительный анализ на конкретной модели DS-CNN с развёртыванием на целевой платформе ESP32-S3 и оценкой влияния на латентность и энергопотребление инференса представляет самостоятельный исследовательский интерес и проводится в настоящей работе.

\subsection{Фреймворки для развёртывания моделей на микроконтроллерах}

Для выполнения инференса нейронных сетей на MCU применяются специализированные среды исполнения (inference runtime), адаптированные к ограничениям встраиваемых платформ: фиксированный объём оперативной и флеш-памяти, отсутствие операционной системы общего назначения и динамической аллокации.

\textit{TensorFlow Lite Micro (TFLM)}~\cite{src:4} является облегчённой версией фреймворка TensorFlow для MCU. Размер библиотеки составляет порядка 100--200~КБ машинного кода. TFLM использует статически выделяемую область памяти (tensor arena) для размещения всех промежуточных тензоров, что исключает необходимость динамической аллокации. Модели хранятся в формате FlatBuffers (.tflite), содержащем граф вычислений, параметры квантизации и массив весов. Фреймворк поддерживает полную INT8-квантизацию входных, выходных и промежуточных тензоров. Для платформ Espressif доступна оптимизированная сборка \texttt{esp-tflite-micro}, интегрирующая библиотеку ESP-NN~--- набор ядер свёрточных операций, реализованных с использованием SIMD-инструкций архитектуры Xtensa LX7~\cite{src:9}.

\textit{ESP-SR (Espressif Speech Recognition)}~\cite{src:14}~--- проприетарная библиотека Espressif для задач распознавания речи, включающая модели WakeNet различных версий. Модели обучены с применением QAT и оптимизированы под архитектуру ESP32-S3, демонстрируя время инференса порядка единиц миллисекунд на фрейм. Основным ограничением является закрытость архитектуры и весов моделей, что исключает возможность модификации, воспроизводимого исследования и сравнительного анализа.

\textit{CMSIS-NN}~\cite{src:13}~--- библиотека оптимизированных нейросетевых ядер для платформ ARM Cortex-M, разработанная Arm Ltd. Реализует свёрточные, полносвязные и pooling-операции с использованием SIMD-инструкций DSP Extension. Является стандартным решением для семейств STM32 и nRF52, однако неприменима к архитектуре Xtensa.

Для настоящей работы выбран TensorFlow Lite Micro с оптимизациями ESP-NN: данный стек обеспечивает развёртывание самостоятельно обученных квантизованных моделей DS-CNN на ESP32-S3 при сохранении аппаратного ускорения и полной воспроизводимости эксперимента.

\subsection{Обоснование выбора целевой платформы}

В качестве платформы для экспериментальной апробации выбран микроконтроллер ESP32-S3 производства Espressif Systems~\cite{src:9}. Выбор обусловлен совокупностью технических и практических факторов.

С точки зрения вычислительной архитектуры, ESP32-S3 построен на двухъядерном процессоре Xtensa LX7 с тактовой частотой до 240~МГц. Согласно техническому справочнику на платформу~\cite{src:18}, архитектура LX7 содержит 128-битные SIMD-регистры (расширение Xtensa PIE), позволяющие выполнять параллельные MAC-операции над данными формата INT8: в один регистр помещаются 16 восьмибитных элементов, что обеспечивает 16 умножений с накоплением за минимальное количество тактов. Данная особенность напрямую ускоряет pointwise-свёртки DS-CNN, являющиеся основной вычислительной нагрузкой модели.

С точки зрения подсистемы памяти, модуль ESP32-S3-Zero N8R8, использованный в настоящей работе, содержит 8~МБ встроенной флеш-памяти и 8~МБ внешней PSRAM. Объём флеш-памяти позволяет разместить прошивку (порядка 1~МБ), библиотеку TFLM (порядка 150~КБ) и квантизованную модель DS-CNN (порядка 260~КБ) с запасом. PSRAM используется для хранения крупных промежуточных буферов.

С точки зрения периферии, ESP32-S3 содержит аппаратный контроллер I2S с поддержкой стандартов Philips, MSB и PCM для подключения цифровых MEMS-микрофонов; контроллеры I2C и SPI для датчиков; модули Wi-Fi~802.11~b/g/n и Bluetooth~5~(LE); таймеры RTC.

С точки зрения режимов энергосбережения, ESP32-S3 поддерживает режим глубокого сна (deep sleep), в котором основные ядра, цифровая периферия, Wi-Fi и Bluetooth обесточены, а активными остаются только домен RTC и сопроцессор ULP. Потребление кристалла в данном режиме составляет порядка 7--8~мкА~\cite{src:9}. Пробуждение инициируется по сигналу на выводах RTC GPIO (механизм EXT0/EXT1), по таймеру RTC или по сигналу от сопроцессора ULP. Выводы RTC GPIO сохраняют заданное логическое состояние в режиме глубокого сна, что позволяет управлять питанием внешней периферии без дополнительных транзисторных ключей.

С точки зрения программной экосистемы, платформа поддерживается открытым фреймворком ESP-IDF~\cite{src:17} (на базе FreeRTOS), содержащим драйверы для всей встроенной периферии, а также компонентом \texttt{esp-tflite-micro} с библиотекой ESP-NN.

С точки зрения доступности, отладочные платы ESP32-S3 доступны на российском рынке по цене порядка нескольких сотен рублей, что существенно ниже стоимости ASIC, FPGA или DSP-решений.

Следует подчеркнуть, что предлагаемая каскадная архитектура не привязана к конкретной платформе. Принципы каскадного пробуждения с возрастающей вычислительной сложностью применимы к любому MCU с аппаратной поддержкой режимов глубокого сна и внешнего пробуждения: STM32U5 с механизмом LPBAM, nRF52840 с функцией wake-on-sound и др. Выбор ESP32-S3 обусловлен задачей экспериментальной апробации и совокупностью практических факторов.

\subsection{Выводы по обзору и постановка задачи}

Проведённый анализ предметной области позволяет сформулировать следующие выводы.

\begin{enumerate}
	\item Задача распознавания ключевых слов на автономных IoT-устройствах характеризуется фундаментальным противоречием между необходимостью непрерывного мониторинга аудиопотока и ограниченным энергетическим бюджетом устройства.
	\item Облачные решения непригодны для батарейных автономных устройств вследствие высокого потребления радиомодуля (600--800~мВт при Wi-Fi), сетевой латентности и зависимости от инфраструктуры.
	\item Специализированные ASIC обеспечивают наивысшую энергоэффективность (единицы мкДж/инференс), но закрыты, дороги и недоступны на российском рынке. FPGA демонстрируют латентность менее 20~мкс и энергопотребление от 30~мкДж/инференс, но требуют специализированных компетенций в проектировании цифровых схем и не обладают IoT-периферией. DSP оптимизированы для классической обработки сигналов, но не для произвольных нейросетевых архитектур.
	\item MCU общего назначения обеспечивают оптимальный баланс гибкости, стоимости, доступности и энергоэффективности при условии применения каскадной архитектуры с управлением режимами энергопотребления.
	\item Каскадная архитектура с этапами возрастающей вычислительной сложности позволяет снизить среднее потребление на порядки за счёт исключения нейросетевого инференса в периоды тишины. Существующие работы рассматривают каскад, квантизацию и декомпозицию энергобюджета по отдельности, но не в совокупности.
	\item Архитектура DS-CNN является эталонной для задачи KWS на MCU, обеспечивая точность свыше 95\% при размере модели 80--500~КБ в зависимости от варианта масштабирования.
	\item Квантизация INT8 обеспечивает сжатие модели до $4\times$ и ускорение инференса за счёт SIMD-инструкций при потере точности менее 1 процентного пункта.
\end{enumerate}

На основании изложенного сформулирована цель настоящей работы: снижение энергопотребления устройств интернета вещей с функцией голосового управления за счёт применения каскадного метода управления режимами работы и использования квантизированных нейронных сетей для распознавания ключевых слов голосовых команд.

Для достижения поставленной цели определены следующие задачи:

\begin{enumerate}
	\item провести анализ существующих методов и систем распознавания ключевых слов, выявить основные источники избыточного энергопотребления и сформулировать требования к энергоэффективному решению;
	\item предложить каскадную архитектуру детекции, разделяющую процесс на этапы с возрастающей вычислительной сложностью и управлением режимами энергопотребления;
	\item исследовать влияние методов квантизации (PTQ, QAT) на точность, размер и вычислительную сложность нейросетевой модели;
	\item провести экспериментальную апробацию на платформе ESP32-S3 с использованием лабораторного стенда и прецизионного измерения энергопотребления;
	\item выполнить сравнительную оценку энергоэффективности предложенного решения и проанализировать вклад каждого компонента каскадной архитектуры в общий энергобюджет.
\end{enumerate}
